\begin{table}[!htbp]
\centering
\caption{Harmonized Peer/Rank Diagnostic Across Experiments}
\label{tab:pooled_peer_effects}
\begin{threeparttable}
\begin{tabular}[t]{lcccc}
\toprule
 & (1) & (2) & (3) & (4) \\
\cmidrule(l{3pt}r{3pt}){2-5}
 & Kenya & Liberia & Nigeria & Pooled \\
\midrule
Peer baseline ability (leave-out mean) &  -0.133** &  -0.048 &  -0.121 &  -0.105** \\
 & (  0.052) & (  0.108) & (  0.098) & (  0.046) \\
\addlinespace
N & 4954 & 3153 & 1659 & 9766 \\
\bottomrule
\end{tabular}
\begin{tablenotes}[para]
\item Notes: This diagnostic applies a harmonized approximate Borusyak--Hull-style specification across the three experiments, with baseline-decile $\times$ treatment $\times$ grade FE and country-nested strata FE. The peer variable is harmonized leave-self-out baseline ability (\texttt{peer\_eb} for Kenya/Liberia, \texttt{peer\_lo\_bl} for Nigeria). SEs are clustered at the harmonized assignment cluster. The pooled column is a diagnostic common-coefficient summary, not a structural peer-effect estimate.
\end{tablenotes}
\end{threeparttable}
\end{table}
